% Abstract (unnumbered chapter)
\chapter*{ABSTRACT}
\addcontentsline{toc}{chapter}{Abstract}

Flashcard apps such as Anki and Duolingo decide when a learner should see each item again. To do this well they need a good
estimate of the chance that the learner still remembers the item today: if the estimate is too high, items are forgotten,
and if it is too low, time is wasted on items the learner already knows. We chose this problem because better estimates
directly mean less study time for the same amount remembered, and because the models in use today are rarely compared on
the same data under the same rules. This report documents and compares several such models on real review histories from
Anki and Duolingo: FSRS, the scheduler that Anki ships; HLR, an older model from Duolingo; our Bayesian memory model, which
tracks a hidden memory strength that fades between reviews, together with two faster approximations of it; and a model
built from many small decision trees (a gradient-boosted tree model) that learns from the learner's own review history. On
the public srs-benchmark leaderboard we also compare against the published neural-network models.

No single model is best at everything, and we report where each one wins. FSRS is a strong and practical baseline: it
keeps a simple memory state that a scheduler can use directly, and on Anki it predicts better than our Bayesian model. The
Bayesian model is the easiest to interpret and was the most accurate of the original models on Duolingo, but its sampler
did not fully converge, so those results are provisional. The tree model is the most accurate model on both of our own
datasets. After we extended it with a richer description of each learner's history and a small per-learner adjustment
that updates after every review, it ranks second of all published models on the public srs-benchmark, ahead of FSRS and
the other neural networks, and behind only the largest neural network (RWKV-Instant). It reaches that position with a
model much smaller than that network, with no GPU, and with predictions roughly a thousand times faster. The tree model
predicts recall but has no memory state of its own, so it complements rather than replaces a scheduler such as FSRS.

\vspace{1em}
\noindent\textbf{Keywords:} spaced repetition; memory modelling; recall prediction; Bayesian models; gradient-boosted
trees; benchmarking.
